% Troubleshooting Guide
% Source: docs/troubleshooting.md (migrated to LaTeX)
% Uses macros from preamble.tex

\section{Troubleshooting Guide}
\label{sec:troubleshooting}

\subsection{Maintaining This Guide}
\label{sec:maintaining-guide}

\textbf{CRITICAL:} Update this file immediately when you encounter and solve new errors. Future work depends on these patterns being documented.

\paragraph{When to add entries:}

\begin{itemize}
  \item \textbf{Always} when you hit a new error that takes more than 10 minutes to debug
  \item When existing error patterns manifest in new ways
  \item After discovering a non-obvious cause for a common symptom
  \item When you find a better solution to an existing problem
  \item After adding new features that create new failure modes
\end{itemize}

\paragraph{How to structure entries:}

\begin{lstlisting}[style=bash]
### Error Title (Brief, Searchable)

**Symptoms:** What the user sees (error messages, wrong output)
**Cause:** Root cause explanation
**Solutions:** Numbered steps to fix
**Don't:** Anti-patterns that seem like they'd work but don't
\end{lstlisting}

\paragraph{Organization:}

\begin{itemize}
  \item Group by subsystem: Wolfram/xAct, Python solver, Pipeline Integration
  \item Keep ``Debugging Techniques'' and ``Verification Checklist'' sections at the end
  \item Cross-reference with \texttt{MEMORY.md} for architectural context
  \item Link to example-specific notes (like \texttt{chern-simons-notes.md}) for complex cases
\end{itemize}

\paragraph{Pruning:} Remove entries if:

\begin{itemize}
  \item The underlying code changed and the error can't happen anymore
  \item Better solutions made the workaround obsolete (note the improvement in \texttt{MEMORY.md})
\end{itemize}


\subsection{Common Wolfram/xAct Issues}
\label{sec:wolfram-xact-issues}

\subsubsection{``Symbol X is already used as a manifold''}

\textbf{Symptoms:} Error when running script multiple times in same kernel session.

\textbf{Cause:} xAct caches tensor definitions in kernel memory.

\textbf{Solutions:}

\begin{enumerate}
  \item Always check before defining:

\begin{lstlisting}[style=mathematica]
If[!xTensorQ[M2], DefManifold[M2, 3, {a,b,c,d,e,f}]]
\end{lstlisting}

  \item Use standard shared symbols across examples:
  \begin{itemize}
    \item Manifolds: \texttt{M2} (1+1D), \texttt{M3} (2+1D), \texttt{M4} (3+1D)
    \item Metrics: \texttt{eta}, \texttt{eta3}, \texttt{eta4}
    \item Covariant derivatives: \texttt{CD}, \texttt{CD3}, \texttt{CD4}
    \item Charts: \texttt{cart}, \texttt{cart3}, \texttt{cart4}
  \end{itemize}

  \item Restart kernel if needed: \texttt{Quit[]} or \texttt{wolframscript -file} (fresh process each time)
\end{enumerate}

\textbf{Don't:} Use \texttt{RandomInteger} or timestamp-based symbol names---kernel caching makes them fail.

\subsubsection{Component Equations Return All Zeros}

\textbf{Symptoms:} After \texttt{DecomposeToComponents}, all equations show \texttt{0}.

\textbf{Likely causes:}

\begin{enumerate}
  \item \textbf{Field strength not expanded:}

\begin{lstlisting}[style=mathematica]
(* BAD *)
eom = VarD[A[-a], CD][L]  (* L contains F[-a,-b] *)
components = DecomposeToComponents[eom, A[-a], cart]  (* F not expanded *)

(* GOOD *)
eom = VarD[A[-a], CD][L /. F[-a,-b] -> CD[-a][A[-b]] - CD[-b][A[-a]]]
(* OR construct L directly in terms of CD[A] *)
\end{lstlisting}

  \item \textbf{Missing field in ToBasis conversion:}
  \begin{itemize}
    \item Check that field symbols appear in \texttt{componentEq} after \texttt{ToBasis[chart][eom]}
    \item If abstract field not converted, check field definition and chart compatibility
  \end{itemize}
\end{enumerate}

\subsubsection{Cross-Field Terms Not Detected}

\textbf{Symptoms:} JSON shows all terms referencing same field, no cross-coupling.

\textbf{Cause:} Other fields not transformed to coordinate form.

\textbf{Example problem:}

\begin{lstlisting}[style=mathematica]
(* Chi appears as cplChi[] not cplChi0[t,x] *)
phiEq = -0.5*cplChi[] - 1.0*cplPhi0[t,x] + Derivative[...][cplPhi0][t,x]
\end{lstlisting}

\textbf{Solution:}

\begin{lstlisting}[style=mathematica]
(* Pass all coupled fields *)
phiComponents = DecomposeToComponents[eomPhi, phi[], cart, {chi[]}]
chiComponents = DecomposeToComponents[eomChi, chi[], cart, {phi[]}]
\end{lstlisting}

\textbf{Verification:}

\begin{itemize}
  \item After decomposition, print equations
  \item All field symbols should have coordinate arguments: \texttt{field0[t,x]} or \texttt{field0[t,x,y]}
  \item No bare field symbols like \texttt{field[]}
\end{itemize}

\subsubsection{JSON Coefficients All Show 1.0}

\textbf{Symptoms:} Correct operator/field but coefficient extraction fails.

\textbf{Cause:} Pattern matching in \texttt{ExtractNumericCoefficient} not finding field symbols.

\textbf{Debug steps:}

\begin{enumerate}
  \item Check \texttt{IdentifyMultiFieldTerm} function head extraction
  \item Verify field names match: \texttt{"phi\_0"} $\to$ \texttt{"phi"} $\to$ \texttt{"Phi"} (case variations)
  \item Print intermediate term structure to see actual symbols
\end{enumerate}

\textbf{Recent fix:} Use \texttt{ToLowerCase} for field base name matching (case-insensitive).

\subsubsection{Multiline Lagrangian Parsed Incorrectly}

\textbf{Symptoms:} Second field's terms evaluate to zero.

\textbf{Cause:} Mathematica multiline without explicit \texttt{+}.

\textbf{Solution:}

\begin{lstlisting}[style=mathematica]
(* BAD *)
L = term1
    term2  (* Treated as separate expression *)

(* GOOD *)
L = term1 +
    term2  (* Explicit continuation *)

(* OR *)
L = (
  term1 +
  term2
)
\end{lstlisting}

\textbf{Example from coupled scalars:} Required \texttt{+} before \texttt{(-1/2 CD[-a][chi[]]...)}.

\subsubsection{Epsilon Tensor Not Evaluating to Numbers}

\textbf{Symptoms:} Epsilon tensor like \texttt{epsiloneta3[\{0,-cart3\},\{1,-cart3\},\{2,cart3\}]} remains in output.

\textbf{Cause:} Pattern mismatch in \texttt{EvaluateEpsilonComponents}.

\textbf{Common issues:}

\begin{enumerate}
  \item Wrong chart name (must match exactly)
  \item Index sign pattern not handled (mixed up/down indices)
  \item Function not being called in pipeline
\end{enumerate}

\textbf{Solution:}

\begin{lstlisting}[style=mathematica]
(* Verify epsilon is being evaluated *)
testExpr = epsiloneta3[{0, -cart3}, {1, -cart3}, {2, cart3}];
result = EvaluateEpsilonComponents[testExpr, cart3];
Print["Result: ", result];  (* Should be a number, not symbolic *)
\end{lstlisting}

\textbf{Debug:} If epsilon remains symbolic, check:

\begin{itemize}
  \item The chart variable matches exactly (e.g., \texttt{cart3} not \texttt{cart})
  \item The epsilon tensor name matches pattern (must contain ``epsilon'')
  \item Both covariant (\texttt{-chart}) and contravariant (\texttt{chart}) cases are handled
\end{itemize}

\subsubsection{Mixed Time-Space Derivatives in 2+1D}

\textbf{Symptoms:} Extra ``laplacian'' terms on cross-fields in JSON.

\textbf{Cause:} \texttt{Derivative[1, 1, 0][f][t,x,y]} ($\partial_t \partial_x f$) classified as second-order.

\textbf{Current status:} Known limitation---the time derivative detection pattern is 1+1D only.

\textbf{Workaround:} For now, the extra terms don't break simulation but may cause numerical artifacts.


\subsubsection{Higher-Curvature / Many-Index Theories Timeout}
\label{sec:higher-curvature-timeout}

\textbf{Symptom:} Derivation takes $>$1 hour or times out for Lagrangians containing curvature-squared ($R^2$, $\tilde{R}^2$) or other higher-curvature terms.

\textbf{Cause:} The legacy abstract \texttt{VarD} + \texttt{TraceBasisDummy} pipeline suffers from combinatorial explosion when the Lagrangian has many contracted abstract indices. For example, $\tilde{R}^2$ in PGT produces $\sim$45 contracted indices, leading to $O(\text{dim}^n)$ scaling ($>$77\,min for 2+1D, infeasible for 3+1D).

\textbf{Resolution:} Component-level Euler--Lagrange (now the default since v0.15.0) bypasses the abstract index machinery entirely by differentiating the component Lagrangian directly. This reduces the same derivation to $\sim$5\,s. See \S\ref{sec:component-el} in the Architecture documentation.

\textbf{Note:} If using TIDAL $<$ v0.15.0, upgrade to benefit from the component E--L pipeline.

\subsubsection{4th-Order-in-Time Equations}
\label{sec:4th-order-eom}

\textbf{Symptom:} JSON specification contains equations with \texttt{d4\_t} or \texttt{d3\_t} on the LHS. Simulation fails because solvers only support $\leq$2nd-order-in-time equations.

\textbf{Cause:} Curvature-squared Lagrangians ($R^2$, $\tilde{R}^2$) produce higher-derivative equations of motion.

\textbf{Resolution (v0.33.0+, current):} Automatic mechanical Ostrogradsky reduction was \emph{removed} in v0.33.0 (commit \texttt{ceb6e63}) because it introduced ghost modes that caused long-time-growth artefacts in simulations.  Higher-derivative theories now require an explicit \texttt{[perturbation]} block in the TOML, with a small parameter declared.  The v6 iterative Parker--Simon scheme (see \cref{sec:perturbative-reduction}) handles such theories by treating the higher-derivative correction as a Duhamel source on the second-order base operator, never propagating ghost modes.  Loading a JSON with \texttt{time\_derivative\_order > 2} but no \texttt{[perturbation]} metadata raises a clear \texttt{ValueError} at JSON-load time.

\textbf{Note:} After Parker--Simon reduction, implicit acceleration cross-references ($\partial_t^2 q_j$ in $q_i$'s equation) are handled by the generalised mass-matrix modal solver. See \cref{sec:modal-mass-matrix} and \cref{sec:troubleshoot-singular-mass}.

\textbf{Architectural barrier (constraint promotion):} For theories where the small parameter promotes an algebraic constraint to a higher-derivative dynamical field (e.g. $b_5\,\tilde{R}^2$ in PGT torsion), the equation-side reduction completes but the Lagrangian-side substitution (LPS, \cref{sec:pr-cb-lps}) aborts; the resulting Hamiltonian retains $\partial_t^2$ residues.  See \cref{sec:pr-constraint-barrier} for the architectural barrier and open research directions.

\subsubsection{IBP Failures for Curvature-Squared Theories}
\label{sec:ibp-curvature-squared}

\textbf{Symptom:} Derivation crashes (segfault, stack overflow, or infinite loop) during the canonical pipeline's integration-by-parts (IBP) step for Lagrangians containing $R^2$ or $\tilde{R}^2$ terms.

\textbf{Causes and fixes} (resolved across v0.18.5--v0.22.4):
\begin{itemize}
  \item \textbf{Detection bug:} \texttt{isPureTimeDerivGE2} checked \texttt{Head[Head[expr]]} but for \texttt{Derivative[2,0,0,0][f][t,x,y,z]}, this returns \texttt{Derivative[2,0,0,0]} not \texttt{Derivative}. Fixed to check \texttt{Head[Head[Head[expr]]]}.
  \item \textbf{Mixed derivatives:} IBP incorrectly flagged mixed first-time+spatial derivatives ($\partial_t \partial_x f$) as needing IBP. Changed to require pure $\partial_t^2$ (time order $\geq 2$). Mixed derivatives like $(\partial_t \partial_z h)^2$ caused unbounded recursion.
  \item \textbf{Stack overflow:} \texttt{FreeQ[lagComp, x\_ /; hasHighTimeDerivative[x]]} tested every subexpression against a condition function, causing C-level stack overflow on $>$200-term Lagrangians. Replaced with structural \texttt{Cases} extraction followed by \texttt{AnyTrue} check.
  \item \textbf{Squared powers:} $(\partial_t^2 f)^2$ terms cannot be IBP'd (would introduce $\partial_t^3$). Now skipped instead of recursively expanded.
\end{itemize}

\subsubsection{Bracket Errors in Generated Wolfram Scripts}
\label{sec:bracket-validation}

\textbf{Symptom:} \texttt{wolframscript} crashes or gives syntax errors after a long derivation run ($>$30 minutes).

\textbf{Cause:} The generated \texttt{.wls} script has unbalanced brackets --- typically from Python f-string quoting bugs where curly braces are not properly escaped.

\textbf{Resolution:} Since v0.22.6, \texttt{validate\_wls\_brackets()} automatically checks bracket balance (\texttt{[]}, \texttt{\{\}}, \texttt{()}) in generated scripts before execution, catching f-string bugs immediately. The validator handles Wolfram string literals, nested comments \texttt{(* ... *)}, and \texttt{\textbackslash[Name]} character escapes.

\subsection{Common Python Solver Issues}
\label{sec:python-solver-issues}

\textbf{Note:} TIDAL's solver layer uses SUNDIALS IDA/CVODE + native numpy operators (see \texttt{tidal/solver/}). The former py-pde-based architecture was replaced in February 2026.

\subsubsection{``Unknown operator: X''}

\textbf{Symptoms:} \texttt{ValueError} when loading or running a simulation.

\textbf{Cause:} JSON spec references an operator not in \texttt{OPERATOR\_REGISTRY} (\texttt{tidal/solver/operators.py}).

\textbf{Available operators:} \texttt{identity}, \texttt{laplacian}, \texttt{laplacian\_x}, \texttt{laplacian\_y}, \texttt{laplacian\_z}, \texttt{gradient\_x}, \texttt{gradient\_y}, \texttt{gradient\_z}, \texttt{cross\_derivative\_xy}, \texttt{cross\_derivative\_xz}, \texttt{cross\_derivative\_yz}, \texttt{biharmonic}, \texttt{first\_derivative\_t}, \texttt{mixed\_T\_S1\_S2\_...}

\textbf{Fix:} Check \texttt{OPERATOR\_REGISTRY} in \texttt{tidal/solver/operators.py}. If the operator is valid but missing, add it there.

\subsubsection{Grid Dimension Mismatch}

\textbf{Symptoms:} Error about field dimensions vs grid dimensions.

\textbf{Cause:} Using \texttt{--grid-shape} with wrong number of dimensions for the JSON spec's spatial dimension.

\textbf{Check:} Match grid shape to the JSON spec's \texttt{spacetime.dimension - 1} spatial dimensions:

\begin{itemize}
  \item 1+1D specs: \texttt{--grid-shape 256} (1 spatial dimension)
  \item 2+1D specs: \texttt{--grid-shape 64,64} (2 spatial dimensions)
  \item 3+1D specs: \texttt{--grid-shape 32,32,32} (3 spatial dimensions)
\end{itemize}

\subsubsection{State Size Mismatch}

\textbf{Symptoms:} Error about unexpected state vector size.

\textbf{Cause:} State has wrong number of slots for the equation system.

\textbf{How slot counts work} (via \texttt{StateLayout} in \texttt{tidal/solver/state.py}):

\begin{itemize}
  \item 2nd-order fields: 2 slots each (field + velocity, e.g., \texttt{phi\_0} + \texttt{v\_phi\_0})
  \item 1st-order fields: 1 slot each (field only)
  \item Constraint fields: 1 slot each (solved algebraically)
  \item Example: 3 vector components (2nd order) $\to$ 6 total slots
\end{itemize}

\subsubsection{Non-Periodic Coefficient Warning / Error}

\textbf{Symptoms:} Warning or \texttt{ValueError} at simulation start:
\texttt{"Position-dependent coefficient `...' has N\% jump at the periodic boundary along x (left=..., right=..., leak\_metric=...)"}

\textbf{Cause:} Position-dependent coefficients (e.g., background magnetic field $B(x)$) must be continuous across periodic boundaries. If they aren't, the integration-by-parts identity fails and causes $O(1)$ energy non-conservation that does \textbf{not} improve with finer grids.

The check uses a \emph{leak metric} = $(\text{rel\_jump}) \times (\text{boundary\_significance})$ to estimate the actual energy leak magnitude. Warnings fire when the metric exceeds 0.01; errors when it exceeds 0.25.

\textbf{Solutions:}

\begin{enumerate}
  \item \textbf{Large domain}: Use a domain large enough that the non-periodic coefficient (e.g., dipolar $B \sim 1/r^3$) is negligible at both boundaries. Terminate the simulation before waves reach the edges.
  \item \textbf{Localised profile}: Use a windowed/Gaussian profile that naturally goes to the same value at both boundaries (e.g., $B(x) = B_\text{peak} \cdot \exp(-(x - x_0)^2 / R^2)$).
  \item \textbf{Non-periodic BCs}: Switch to Dirichlet or Neumann BCs (\texttt{--bc neumann}).
  \item \textbf{False positive?} If the warning fires but energy conservation is actually fine (check with \texttt{tidal measure --what conservation}), the coefficient boundary values are small enough that the leak is negligible. The leak metric should be very small in this case---if it isn't, the warning is genuine.
\end{enumerate}

\subsubsection{Constraint Solver Failures}

\textbf{Symptoms:} IDA fails to converge, or constraint fields have NaN values.

\textbf{Cause:} Algebraic constraint equations may have singular operators (e.g., Laplacian with periodic BCs has a null space for constant functions).

\textbf{Solutions:}

\begin{enumerate}
  \item The three-tier constraint pre-solve (\texttt{tidal/solver/constraint\_solve.py}) handles most cases automatically
  \item For periodic BCs with pure-Laplacian constraints, gauge regularisation pins the zero mode
  \item Check \texttt{--ic} settings---constraint ICs must be consistent with dynamical field ICs
\end{enumerate}


\subsubsection{ZeroDivisionError: \texttt{float division by zero} in coefficient evaluation}
\label{sec:troubleshooting-zerodiv-kinetic}

\textbf{Symptom:}
\begin{lstlisting}[style=python]
Cannot evaluate symbolic coefficient '(2*alpha)/xi': float division by zero
\end{lstlisting}
when running \texttt{tidal simulate} or \texttt{tidal sweep} with a parameter set to
zero (commonly \texttt{xi=0} in the dark photon torsion model).

\textbf{Root cause:} Old JSON files (pre-v0.24.3) had \texttt{ExportJSON.wl} divide
the RHS symbolically by the kinetic coefficient at derivation time.  When that
coefficient evaluates to zero the stored expression (e.g.\ \texttt{(2*alpha)/xi}) is
undefined.

\textbf{Fix:} Re-derive the theory with the current pipeline
(\texttt{tidal derive theory.toml}).  Since v0.24.3, \texttt{ExportJSON.wl} detects
parameter-based kinetic coefficients and emits them in the LHS as
\texttt{kinetic\_coefficient\_symbolic} rather than dividing.  The Python function
\texttt{normalize\_kinetic\_coefficients(spec, params)} in
\texttt{tidal.symbolic.json\_loader} then handles normalisation at runtime:
\begin{itemize}
  \item $K \neq 0$: divides through correctly.
  \item $K = 0$: transitions the field to a constraint (\texttt{time\_order = 0}).
\end{itemize}

\textbf{Diagnostic:} Check whether the offending equation's LHS has
\texttt{kinetic\_coefficient\_symbolic} in the JSON:
\begin{lstlisting}[style=python]
import json
d = json.load(open("examples/data/my_theory.json"))
for eq in d["equations"]:
    kc = eq.get("lhs", {}).get("kinetic_coefficient_symbolic")
    if kc:
        print(eq["field"], "->", kc)
\end{lstlisting}
If the field is absent, the JSON pre-dates v0.24.3 --- re-derive.

\textbf{Physics note:} Setting $\xi = 0$ in the dark photon torsion model removes the
torsion kinetic term; the torsion fields become algebraic constraints and the system
reduces to the standard Einstein-Maxwell Gertsenshtein model.  This is a useful
control check: the conversion probability should recover $P = \sin^2(\kappa B_0 D/2)$.


\subsection{Debugging Techniques}
\label{sec:debugging-techniques}

\subsubsection{Wolfram Side}

\begin{enumerate}
  \item \textbf{Print intermediate steps:}

\begin{lstlisting}[style=mathematica]
Print["After VarD: ", eom];
Print["After ToBasis: ", ToBasis[chart][eom]];
Print["After metric eval: ", EvaluateMinkowskiMetric[expr, chart]];
\end{lstlisting}

  \item \textbf{Check field transformation:}

\begin{lstlisting}[style=mathematica]
(* Should see coordSyms = {t[], x[]} or {t[], x[], y[]} *)
coordSyms = GetCoordinateSymbols[chart]
Print[coordSyms]
\end{lstlisting}

  \item \textbf{Verify dimension:}

\begin{lstlisting}[style=mathematica]
dim = Length[ScalarsOfChart[chart]]
Print["Dimension: ", dim]
\end{lstlisting}

  \item \textbf{Test ToCanonical:}

\begin{lstlisting}[style=mathematica]
(* Simplify before and after to see if indices contract properly *)
Print["Before: ", expr];
expr = ToCanonical[expr];
expr = ContractMetric[expr];
Print["After: ", expr];
\end{lstlisting}
\end{enumerate}

\subsubsection{Python Side}

\begin{enumerate}
  \item \textbf{Validate JSON load:}

\begin{lstlisting}[style=python]
spec = load_equation_system(json_path)
print(f"Dimension: {spec.dimension}")
print(f"Components: {spec.component_names}")
for eq in spec.equations:
    print(f"{eq.field_name}: {len(eq.rhs_terms)} terms")
\end{lstlisting}

  \item \textbf{Check state structure:}

\begin{lstlisting}[style=python]
print(f"State fields: {len(state)}")
for i, field in enumerate(state):
    print(f"  {i}: {field.label}, shape={field.data.shape}")
\end{lstlisting}

  \item \textbf{Monitor evolution:}

\begin{lstlisting}[style=python]
# Add print inside evolution_rate or use callback tracker
def check_rates(state):
    print(f"Max field value: {max(np.max(np.abs(f.data)) for f in state)}")
\end{lstlisting}
\end{enumerate}


\subsection{Verification Checklist}
\label{sec:verification-checklist}

\subsubsection{After Wolfram Changes}

\begin{itemize}
  \item Run all examples: \texttt{./scripts/run\_examples.sh} (or \texttt{tidal derive examples/*/theory.toml})
  \item Check JSON dimension matches spacetime (2 for 1+1D, 3 for 2+1D)
  \item Verify component count matches field rank $\times$ dimensions
  \item Spot-check coefficients in JSON (shouldn't all be 1.0)
  \item Confirm cross-field terms if applicable
\end{itemize}

\subsubsection{After Python Changes}

\begin{itemize}
  \item Run pytest: \texttt{uv run pytest tests/}
  \item Test JSON loading for all examples
  \item Run at least one simulation end-to-end
  \item Check output plots if visualization enabled
  \item Verify energy conservation (if applicable to physics)
\end{itemize}

\subsubsection{After Pipeline Changes}

\begin{itemize}
  \item Regression test: \texttt{coupled\_scalars} still works
  \item Forward test: new example runs
  \item Cross-test: modify old example to use new feature
  \item Documentation: update \texttt{MEMORY.md} with new patterns
\end{itemize}


\subsection{Higher-Derivative Theory Issues}
\label{sec:troubleshoot-higher-derivative}

\subsubsection{Singular Mass Matrix}
\label{sec:troubleshoot-singular-mass}

\textbf{Symptom}: Modal solver logs ``Generalized mass matrix: singular M detected'' and produces large matrix entries ($\max|A| \gg 1$).

\textbf{Cause}: A singular mass matrix typically indicates hidden algebraic constraints between fields, which can arise from gauge symmetry (e.g., Poincar\'e gauge symmetry in PGT) or from constraint-field demotion under v6 perturbative reduction (\cref{sec:pr-v6}). Pre-v6 derivations also produced singular $M$ via Ostrogradsky-auxiliary fields, but the Ostrogradsky reduction was removed in v0.33.0; see \cref{sec:4th-order-eom}.

\textbf{Solution}: The modal solver handles this automatically via mass-matrix eigendecomposition (\cref{sec:modal-singular-mass}). Zero eigenvalues are detected and the corresponding directions are Schur-eliminated. No user action needed.

\textbf{Example}: In graviton-torsion (PGT), $M = \bigl[\begin{smallmatrix} 1 & 1 \\ 1 & 1 \end{smallmatrix}\bigr]$ for $\{t_3, t_7\}$ reveals the constraint $t_3 = t_7$ from Poincar\'e gauge symmetry.

\subsubsection{Velocity Coupling Singularity}
\label{sec:troubleshoot-velocity-coupling}

\textbf{Symptom}: Modal solver logs ``Generalized eigenvalue: N gauge modes zeroed'' during simulation.

\textbf{Cause}: Constraint fields with circular velocity dependencies (e.g., $h_0$ references $v_{h_1}$ while $h_1$ references $v_{h_0}$) create a singular implicit velocity equation $(I - V_\text{coupling}) \cdot d' = A \cdot d$.

\textbf{Solution}: Resolved automatically via QZ generalized eigenvalue decomposition (\cref{sec:modal-qz-velocity}). Infinite eigenvalues correspond to gauge degrees of freedom and are frozen at initial values.

\textbf{History}: Prior to v0.16.1, this was handled by regularization ($+\epsilon I$), which produced spurious $O(10^8)$ oscillation frequencies. The QZ fix (commit \texttt{e2e9cc7}) eliminates these entirely.

\subsubsection{Rank-Deficient Kinetic from Trace-Projection Lagrangians}
\label{sec:troubleshooting-rank-deficient-kinetic}

\textbf{Symptom}: Simulation of a PGT-style model diverges with ``\texttt{Simulation diverged at t=\ldots: amplitude ratio \ldots exceeds threshold 1e+08}'' and the per-mode warning ``\texttt{max Re($\lambda$)$\approx$1}'' even at physically benign parameters. Typically surfaces for theories that linearize a higher-rank tensor (e.g. the 24-component PGT torsion \texttt{TorsionCDT[a,-b,-c]} perturbation) but only use a low-rank projection (e.g. the trace vector $T_a = T^b{}_{ba}$) in the Lagrangian.

\textbf{Cause}: After xPert linearization, the JSON contains all components of the rank-3 perturbation, but the Lagrangian terms $\xi\,\mathrm{Ftorsion}^2$ and $\alpha\,I_3$ only place kinetic and mass on the 4-dim trace subspace. The resulting dynamical block has a \emph{rank-deficient mass matrix} $M$ --- for example a 3-component trace triple $(t_2, t_{10}, t_{17})$ gives $M = \bigl[\begin{smallmatrix} 1 & 1 & 1 \\ 1 & 1 & 1 \\ 1 & 1 & 1 \end{smallmatrix}\bigr]$ with eigenvalues $(3, 0, 0)$ --- two null directions that represent gauge/non-propagating degrees of freedom rather than physical modes. Prior to \#256 the modal solver routed such theories through a separate constraint-elimination builder (\texttt{build\_constraint\_eliminated\_matrices}) that silently mis-classified \texttt{d2\_t} cross-couplings as stiffness operators (uses \texttt{\_EXACT\_MULTIPLIERS}, which returns only the spatial part of the operator; for \texttt{d2\_t} this is $\mathtt{ones\_like}(k)=1$), producing spurious tachyonic modes with $\mathrm{Re}(\lambda) \approx 1$ that blew up over long evolutions.

\textbf{Resolution}: The unified builder \texttt{\_build\_evolution\_matrices} detects singular $M$ at line~$\sim$860 of \texttt{tidal/solver/modal.py} via \texttt{np.linalg.det(M) < 1e-12}, eigendecomposes $M$, partitions into dynamical ($\lambda \neq 0$) and mass-constrained ($\lambda = 0$) blocks, and Schur-eliminates the null subspace. Model authors see a benign INFO log line:
\begin{verbatim}
Generalized mass matrix: singular M detected — applying
mass-matrix Schur elimination
Mass matrix eigenvalues: [0, 0, ..., 3, 3, 3, ...]
(dynamical: N, constrained: M)
\end{verbatim}
This is normal and expected --- no user action needed.

\textbf{Example}: The \texttt{torsion\_dark\_photon} model with Lagrangian $(1/\kappa^2) R + \alpha I_3 - \frac{1}{4}\xi\,\mathrm{Ftorsion}^2 + \delta_m\,F\cdot\mathrm{Ftorsion} - \frac{1}{4}F^2$ at $(\alpha{=}0.5,\xi{=}1,\delta_m{=}0.1)$ used to diverge at $t \approx 34$ with amplitude ratio $10^{8}$. After the \#256 unification it runs to $t_\mathrm{end}=50$ producing $P_\mathrm{max}(h_5\to a_1) = 0.061197$, matching $\sin^2(\kappa B_0 t/2) = \sin^2(0.25) = 0.061209$ to four digits. Ratio $P(50)/P(25) = 4.05$ confirms $\sin^2$ growth in the linearized regime.

\textbf{Related}: \#256 (solver unification), \#166 (near-singular velocity coupling), \#175 (massive-gravity energy conservation), \#222 (\texttt{\_suppress\_tachyonic\_noise} for zero-coupling tachyons). See also \cref{sec:troubleshoot-singular-mass} above for the singular-$M$ path that \#256 unified.

\subsubsection{Ghost Instability Warning}
\label{sec:troubleshoot-ghost}

\textbf{Symptom}: Modal solver warns ``eigenvalues with positive real parts (max Re($\lambda$)=...)'' and energy grows exponentially.

\textbf{Cause}: The Ostrogradsky theorem states that non-degenerate Lagrangians with time derivatives $> 1$ generically propagate ghosts (unbounded Hamiltonian). Some PGT parameter windows are ghost-free (Barker 2024, arXiv:2406.12826).

\textbf{Solution}: This is physical, not a bug. Check:
\begin{itemize}
  \item Whether the theory parameters are in a ghost-free window
  \item Whether the growth rate is physical or numerical (compare with IDA/CVODE)
  \item The $(-,+,+,+)$ metric convention: negative kinetic energy for timelike components is physical, not a ghost
\end{itemize}
See issue \#164 for full propagator analysis (future work).
